\documentclass[a4paper,12pt]{article}
% \documentclass[fleqn,12pt]{article}
\usepackage[a4paper,top=2.5cm,bottom=2.5cm,left=2cm,right=2cm]{geometry}
\usepackage{amsmath}
\usepackage{times}
\usepackage{graphicx}
\usepackage[english,danish]{babel}
\usepackage[utf8]{inputenc}
\usepackage[T1]{fontenc}
\usepackage{booktabs}
\usepackage{hyperref}
\usepackage{float} 
\usepackage{minted}
\usepackage{caption}  % for caption* to include notes under the table
% Bibliography style
\usepackage{natbib}        % For citations and references
\bibliographystyle{econometrica}


% set line spacing
\renewcommand{\baselinestretch}{1.2}
\title{Assignment 3, Økonometri A}
\author{}
% \date{1. november, 2024, kl. 10.00-22.00}

\begin{document}

\maketitle

\section*{Introduktion til opgaven}
En afgørende udfordring i analyser af bilmarkedet er at forstå, hvordan forbrugerne reagerer på ændringer i priser og produktkarakteristika, og hvordan virksomheder fastsætter deres priser under hensyntagen til disse reaktioner. Ved at kombinere data med økonomisk teori og statistiske metoder kan vi opnå en dybere indsigt i disse mekanismer og kvantificere effekterne af forskellige politikker på bilmarkedet.

Vi tager udgangspunkt i en strukturel model for udbud og efterspørgsel efter biler inspireret af \cite{grieco2024evolution}. Vi undersøger først, hvordan forbrugere træffer valg mellem forskellige bilmodeller baseret på bilernes priser og egenskaber. Under antagelse om optimal prissætning undersøger vi herefter, hvordan vi kan udlede virksomhedernes (uobserverede) marginalomkostninger. Denne strukturelle tilgang lader os estimere parametre, der kan bruges til at simulere effekterne af forskellige politikker på bilmarkedet, såsom skatter på hestekræfter eller ændrede tariffer.

Afleveringen rummer fire opgaver, som hver består af tre delopgaver. Opgave 1-3 fokuserer på empiriske analyser af bilmarkedet, mens opgave 4 omfatter en teoretisk og  simulationsbaseret Monte Carlo-analyse. Opgave 4 kan løses uafhængigt af opgave 1-3. Nogle studerende vil have fordel af at starte med opgave 4, da arbejdet med teorien kan hjælpe med fortolkningen af de empiriske resultater.
\\

\noindent \textbf{Bemærk: Din besvarelse må maksimalt fylde 2600 ord (dvs. ca. 10 sider). Antallet af ord skal fremgå på din forside}. Hvis du skriver din opgave i Overleaf, kan du tælle opgavens ord ved at klikke Menu $\rightarrow$ Word Count øverst til venstre. Figurer, tabeller og forside tæller ikke med i dit sideantal. 

Det anbefales at du er kortfattet, præcis og formel i din argumentation, og at du konkret henviser til modellen og den økonometriske teori, når det er muligt.
\subsection*{Data}
Til denne aflevering medfølger datasættet \texttt{cars.csv}, som er en modificeret udgave af datagrundlaget for \cite{grieco2024evolution}. Indlæses filen som DataFrame i Python med koden:

\begin{minted}[fontsize=\footnotesize, baselinestretch=1.0, breaklines, frame=single]{python}
    import pandas as pd
    df = pd.read_csv("cars.csv") # Tilpas evt. stien til csv-filen
    df.describe(percentiles=[]).T
    \end{minted}
\noindent
\noindent Datasættet indeholder 9.694 observationer af markedsandele, priser og produktkarakteristika for forskellige bilmodeller på det amerikanske marked fordelt på årene 1980-2018. Hvert produkt $j$ er identificeret ved en bilproducent (\texttt{make}) og en bilmodel (\texttt{model}) i år $t$ (\texttt{year}). Observationsenheden er altså produkt-år kombinationer $(j,t)$, og et eksempel er: \texttt{Toyota Corolla, 2007}.

Priser (\texttt{price}) er målt i \$1.000 (2015-priser). \texttt{year} og \texttt{number\_households} varierer kun over tid. Dummyvariablene \texttt{car}, \texttt{suv}, \texttt{truck}, og \texttt{van} er gensidigt udelukkende og summerer til 1. Dummyvariablene \texttt{EV}, \texttt{HEV}, og \texttt{PHEV} er også gensidigt udelukkende, men de summerer ikke til 1 (f.eks. mangler dummy for diesel- og benzinbiler). De resterende to dummier, \texttt{sport} og \texttt{lux}, er ikke gensidigt udelukkende og der er ingen dummy for referencekategorien. 

Tabel \ref{tab:variable_description} giver et overblik over alle de tilgængelige variabler.

\begin{table}[H]
    \centering
    \caption{Variabelbeskrivelse}
    \footnotesize
    \begin{tabular}{p{3cm}p{11cm}}
    \toprule
    \textbf{Variabel} & \textbf{Beskrivelse og Værdisæt} \\
    \midrule
    {year} & Kalenderår for observation. Interval, $t=1980,1981 \dots 2018$ \\
    {product\_id} & Unik ID for make-model,  $j=0,\dots,J=1066$ \\
    {make} & Bilproducent som Ford, Toyota, BMW, Volvo, etc. \\
    {model} & Model af køretøj som BMW 3 Series, Chevrolet Suburban, Ford Mustang, etc. \\
    {number\_households} & Antal husholdninger i året, $n_t$ \\
    {sales} & Solgte enheder i et givent år \\
    {price} & Anbefalet udsalgspris, $p_{jt}$, målt i \$1,000 (2015-priser)  \\
    {height} & Højde på køretøjet. Målt i tommer \\
    {footprint} & Areal (længde × bredde). Målt i kvadrattommer \\
    {weight} & Vægt af køretøjet. Målt i pund \\
    {hp} & Hestekræfter. \\
    {mpg} & Brændstofeffektivitet målt i miles per gallon \\
    {yearsSinceDesign} & Antal år siden design blev introduceret. Interval: 0-25\\
    {number\_trims} & Antal trimniveauer for modellen (dvs. konfigurationsmuligheder) \\
    {releaseYear} & Dummy for nyt modelår \\
    {lux} & Dummy for luksusbiler \\
    {sport} & Dummy for sportsbiler \\
    {HEV} & Dummy for hybridbiler \\
    {PHEV} & Dummy for plug-in hybridbiler \\
    {EV} & Dummy for elbiler \\
    {car} & Dummy for personbiler \\
    {suv} & Dummy for SUV'er \\
    {truck} & Dummy for lastbiler \\
    {van} & Dummy for varevogne \\
    {market\_share} & Markedsandel, $s_{jt} \in (0,1)$ \\
    {no\_car} & Andel af husholdninger uden bil, $s_{0t} \in (0,1)$ \\
    {dlogS} & Logaritme af markedsandele relativt til andelen uden bil,  
    $\log \left(\frac {s_{jt}} {s_{0t}} \right)$ \\
    {RXR} & Real valutakurs \\
    {gh\_log\_$g$} & GH-instrumentet for log $g \in$ \{height, footprint, hp, mpg, weight\} \\
    \bottomrule
    \\
    \end{tabular}
    \label{tab:variable_description}
    \caption*{Note: Nogle variabler er udeladt, omdøbt, skaleret anderledes eller tilføjet ift. det oprindelige datasæt. }
\end{table}

\newpage

%%%%%%%%%%%%%%%%%%%%%%%%%%%%%%%%%%%%%%%%%%%%%%%%%%%%%%%%%%%%%%%
% Q1: Descriptive Analysis and OLS Estimation
%%%%%%%%%%%%%%%%%%%%%%%%%%%%%%%%%%%%%%%%%%%%%%%%%%%%%%%%%%%%%%%
\section*{Model (opsummering)}
Afleveringen tager udgangspunkt i en strukturel model for bilmarkedet, som kan opsummeres med følgende tre lineære ligninger:

\vspace{1em}
\noindent \emph{1. Markedsandele (efterspørgselskurve)}
\begin{equation}
    \log\left(\frac{s_{jt}}{s_{0t}}\right) = \beta \mathbf{x}_{jt} + \alpha p_{jt} + \xi_{jt}. \label{eq:demand}    
\end{equation}


\noindent \emph{2. Prissætning (udbudskurve)}
\begin{equation}
    p_{jt}  = mc_{jt} +\frac{1}{-\alpha (1 - s_{jt})}. \label{eq:supply}    
\end{equation}


\noindent \emph{3. Marginalomkostninger}
\begin{equation}
    \log(mc_{jt}) = \gamma_x \mathbf{x}_{jt} + \gamma_w \mathbf{w}_{jt} + \nu_{jt}, \label{eq:mc}    
\end{equation}
hvor
\begin{itemize}
\itemsep-.5em 
    \item $s_{jt}$ er markedsandelen for produkt $j$ i år $t$. $s_{0t}$ er andelen af forbrugere uden bil i år $t$.
    \item  $\mathbf{x}_{jt}$ er de observerbare produktkarakteristika som også påvirker nytten (f.eks. størrelse).
    \item $p_{jt}$ er den observerbare pris.
    \item  $\xi_{jt}$ er fejlleddet i ligning 1, dvs. uobserverede produktkarakteristika som påvirker nytten.
    \item $mc_{jt}$ er de uobserverede marginalomkostninger.
    \item $\mathbf{w}_{jt}$ er udbudsspecifikke karakteristika (f.eks. råvarepriser, skatter, lønomkostninger, mv.).
    \item $\nu_{jt}$ er fejlledet i ligning 3, som også kan tolkes som
    et stokastisk udbudsstød.
    \item $\beta$, $\alpha$, $\gamma_x$ og $\gamma_w$ er parametre, som vi ønsker at estimere.
\end{itemize}
Bemærk, at ovenstående kun er en kortfattet opsummering af modellen. En komplet beskrivelse af modellen kan findes i appendikset. Det anbefales, at man sætter sig ind i den fulde strukturelle model, så man kan relatere sin analyse til den økonomiske teori.

\vspace{2em}

\noindent\emph{Ligevægt på bilmarkedet}\\
Den centrale udfordring er, at de tre ligninger ovenfor introducerer en simultanitet mellem priser $p_{jt}$ og markedsandele $s_{jt}$. Den simultanitet er vi nødt til at tage højde for, hvis vi ønsker at estimere modellens parametere med data på observerede priser $p_{jt}$, markedsandele $s_{jt}$, produktkarakteristika $\mathbf{x}_{jt}$ og udbudsspecifikke variable $\mathbf{w}_{jt}$. Det virker umiddelbart som en vanskelig opgave, da både markedsandelene $s_{jt}$ og priserne $p_{jt}$ er endogene, fordi de er bestemt i en ligevægt.  Heldigvis tager vi et skridt ad gangen. 

\break
\section*{Opgave 1: Indledende analyse (25\%) }
\noindent \textbf{(1A)} Udfør en deskriptiv analyse af datasættet \texttt{data.csv}. 

\begin{itemize}
    \item Illustrér udviklingen over tid i priser, markedsandele og andre centrale attributter som hestekræfter (\texttt{hp}), brændstofeffektivitet (\texttt{mpg}), areal (\texttt{footprint}) og vægt (\texttt{weight}) for forskellige biltyper (f.eks. \texttt{car}, \texttt{suv}, \texttt{truck}, \texttt{van}).
    \item Undersøg sammenhængen mellem logaritmen af markedsandele relativt til andelen uden bil (\texttt{dlogS}) og bilpriser (\texttt{price}), samt mellem priser og fysiske attributter for bilerne. Kommentér kort på de observerede tendenser.
\end{itemize} 
\hrulefill
\bigskip

\noindent \textbf{(1B)} Brug OLS til at estimere varianter af efterspørgselsligningen i ligning \eqref{eq:demand} som specificeret herunder. Rapporter for hver af specifikationerne dine parameterestimater, standardfejl, TSS, ESS, RSS og $R^2$.
I alle specifikationerne er den afhængige variabel logaritmen af markedsandelen relativt til andelen uden bil, dvs. $y_{jt}= \log\left(\frac{s_{jt}}{s_{0t}}\right)$ (variablen \texttt{dlogS}). Vi ser i første omgang bort fra endogenitetsproblemet, og antager derfor her, at prisen er eksogen.

\begin{itemize}
    \item \textit{Specifikation 1:} Estimer en variant uden produktkarakteristika (men med et konstantled $\beta_0$): 
$$y_{jt}= \beta_0 + \alpha p_{jt} + \xi_{jt}$$
Fortolk fortegnet til $\hat \alpha$. Stemmer det overens med din forventning og almindelige økonomiske antagelser?

\item \textit{Specifikation 2:} Estimer en variant med produktkarakteristika:
$$y_{jt}= \beta \mathbf{x}_{jt} + \alpha p_{jt} + \xi_{jt}$$
hvor matricen $\mathbf{x}_{jt}$ indeholder følgende karakteristika: et konstantled, variablerne  \texttt{hp}, \texttt{mpg}, \texttt{height}, \texttt{footprint}, \texttt{weight}, \texttt{number\_trims}, \texttt{releaseYear}, og \texttt{yearsSinceDesign}, samt dummies for biltyperne \texttt{sport}, \texttt{EV}, \texttt{truck}, \texttt{suv} og \texttt{van}. 

\item \textit{Specifikation 3:} Estimer en variant, hvor vi også kontrollerer for tidseffekten $\delta_t$ ved at inkludere et fuldt sæt af tids-dummies for kalenderår dannet ud fra variablen \texttt{year}:
$$
y_{jt}= \beta \mathbf{x}_{jt} + \alpha p_{jt} + \delta_t + \xi_{jt}
$$
\textit{Hint:} Tidsdummies kan genereres i pandas ved at bruge \texttt{pd.get\_dummies()}. Eksempelvis:


\begin{minted}[fontsize=\scriptsize, baselinestretch=1.0, breaklines, frame=single]{python}
    # Add dummies for 'year'
    year_dummies = pd.get_dummies(df['year'], prefix='y', drop_first=True, dtype=int)
    df = pd.concat([df, year_dummies], axis=1)

    # Store year variable names for later use
    years = list(year_dummies.columns)  # ['y_1981', 'y_1982', ...] 
\end{minted}


\item \textit{Specifikation 4:} Estimer samme specifikation, men hvor du log-transformerer de kontinuerte forklarende variable
\texttt{hp}, \texttt{mpg}, \texttt{height}, \texttt{footprint}, \texttt{weight} og \texttt{number\_trims}.

\item Sammenlign de fire specifikationer og fortolk dine resultater. Diskuter, hvordan kontrol for forskellige variable påvirker $\hat \alpha$. Vurder, hvilken model der bedst beskriver data.
\end{itemize}
\hrulefill
\bigskip

\noindent \textbf{(1C)}
Udfør et test af $H_{0}:\alpha=0$ mod det tosidede alternativ $H_{A}:\alpha\neq0$. 
\begin{itemize}
    \item Opstil den relevante test statistik og redegør for de nødvendige antagelser bag testet.
    \item Konkluder på testresultatet og diskuter, hvordan testresultatet påvirkes af, om der inkluderes produktkarakteristika, biltyper og tidseffekter i modellen.
    \item Diskuter om resultatet er følsomt over for brugen af heteroskedasticitetsrobuste standardfejl.
\end{itemize}

%%%%%%%%%%%%%%%%%%%%%%%%%%%%%%%%%%%%%%%%%%%%%%%%%%%%%%%%%%%%%%%
% Fixed Effects, Frisch-Waugh Theorem, and IV Estimation
%%%%%%%%%%%%%%%%%%%%%%%%%%%%%%%%%%%%%%%%%%%%%%%%%%%%%%%%%%%%%%%
\bigskip
\section*{Opgave 2: Fixed Effects og IV-estimation (25\%)}
\noindent \textbf{(2A)} Vi vil nu udvide vores model ved at inkludere bilmærke-specifikke fixed effects baseret på variablen \texttt{make}. Vi dekomponerer altså efterspørgsels-støddet som \( \xi_{jt} = \bar{\xi}_m + \tilde{\xi}_{jt} \), hvor \( \bar{\xi}_m \) er bilmærke-specifikke effekter, som er konstante over tid for hvert bilmærke, og \( \tilde{\xi}_{jt} \) er den resterende del af fejlledet, der varierer både med bilmodel og tid.

\begin{itemize}
    \item Estimér først modellen ved at inkludere make-specifikke dummies som beskrevet ovenfor, dvs. 
    \begin{equation}
        y_{jt} = \beta \mathbf{x}_{jt} + \alpha p_{jt} + \delta_{t} + \bar{\xi}_m + \tilde{\xi}_{jt}, \label{eq:make_fixed_effects_model}
    \end{equation}
    Inkluder fortsat alle kontrolvariablerne fra specifikation 4 i opgave 1B.

    \textit{Hint}: $\bar{\xi}_m$ kan estimeres ved at inkludere en dummy, som antager værdien 1 hvis bilmærket er $m$ og 0 ellers. Variablen kan genereres  med samme fremgangsmåde, som vi brugte til vores tidsdummies i opgave 1B.


    \item Estimer derefter modellen ved hjælp af within-transformationen:
\begin{equation}
    y_{jt} - \overline{y}_{m} = \beta \left(\mathbf x_{jt} - \overline{\mathbf x}_{m}\right) + \alpha \left(p_{jt} - \overline{p}_{m}\right) + (\delta_{t} - \overline{\delta}_m) + (\tilde{\xi}_{jt} - \overline{\tilde{\xi}}_{m}),
\end{equation}
hvor \( \overline{\cdot}_m \) angiver gennemsnittet inden for hvert bilmærke \( m \) baseret på variablen \texttt{make}.
For eksempel beregnes $\overline{\mathbf x}_{m} = \frac{1}{T_{m}} \sum_{j \in m} \mathbf x_{jt}$, hvor $T_{m}$ er antallet af observationer for bilmærke $ m $. De øvrige gennemsnit, såsom $\overline{p}_{m}$, er defineret på tilsvarende måde. Within-transformationen kan implementeres i pandas ved at bruge \texttt{groupby()} og \texttt{transform()}. Eksempelvis:

\begin{minted}[fontsize=\scriptsize, baselinestretch=1.0, breaklines, frame=single]{python}
df_within = df.copy()
vars_to_transform = [...] # List the names of the variables to transform

for var in vars_to_transform:
    make_mean = df.groupby('make')[var].transform('mean')
    df_within[var] = df[var] - make_mean
\end{minted}

\item Sammenlign parameterestimaterne for de to modeller. Er de ens? Forklar, hvorfor konstantleddet er nul i modellen med within-transformerede variable, og hvorfor \( \bar{\xi}_m \) forsvinder efter within-transformationen (\textit{Hint:} læg mærke til forskellen mellem $\overline{\tilde{\xi}}_{m}$ og $\overline{{\xi}}_{m}$).

\item Sammenlign også TSS, ESS, RSS og \( R^2 \) for de to modeller. Diskuter betydningen af eventuelle forskelle i goodness-of-fit mål. Forklar, hvorfor \( R^2 \) er højere i modellen med make-dummies sammenlignet med den within-transformerede model.

\end{itemize}
\hrulefill

\bigskip 

\noindent \textbf{(2B)} Forklar nødvendigheden af at inkludere de make-specifikke fixed effects $\overline \xi_m$. Tag udgangspunkt i efterspørgselsligningen, hvor vi har dekomponeret fejlledet som beskrevet i 2A:
    $$
    \log\left(\frac{s_{jt}}{s_{0t}}\right) = \beta \mathbf{x}_{jt} + \alpha p_{jt} + \delta_t + \overline \xi_m +  \tilde \xi_{jt}.
    $$

\begin{itemize}
    \item Hvis de uobserverede bilmærke-specifikke effekter $\overline \xi_m$ (fx oplevet brand-image, kvalitetsfølelse) er positivt korrelerede med prisen $p_{jt}$, hvilken retning forventer du så, at bias i estimatet af $\alpha$ vil tage, hvis $\overline \xi_m$ udelades fra modellen?
    \item Er det konsistent med det, du finder? (\textit{Hint:} Sammenlign dine resultater fra 1B med 2A)
\end{itemize}
\hrulefill
\bigskip 

\noindent \textbf{(2C)} Den resterende tids- og bilmodelvarierende komponent \( \tilde{\xi}_{jt} \) kan stadig være korreleret med prisen \( p_{jt} \). For at adressere dette anvender \cite{grieco2024evolution} den reale valutakurs, \texttt{RXR}, som instrument. \texttt{RXR} påvirker især udenlandske bilproducenters marginalomkostninger og indgår derfor i de udbudsspecifikke faktorer \( \mathbf{w}_{jt} \) i ligning \eqref{eq:mc}, som igen påvirker udbudsprisen \( p_{jt} \) gennem ligning \eqref{eq:supply}.

Bemærk, at selv hvis den reducerede form for prisfunktionen \( p_{jt} \) ikke er lineær i de eksogene variable, så er en korrekt specificering af denne \textit{ikke} et krav for at benytte \texttt{RXR} som instrument. Vi kan derfor anvende følgende lineære approksimation af den reducerede form for prisfunktionen:
\begin{equation}
    p_{jt} = \gamma_0 + \gamma_1 \text{RXR}_{jt} + \gamma_2 \mathbf{x}_{jt} + \kappa_{t} + \phi_{m} + v_{jt}, \label{eq:first stage rxr}
\end{equation}
hvor \( \mathbf{x}_{jt} \) er de eksogene produktkarakteristika fra efterspørgselsligningen, \( \kappa_t \) og \( \phi_m \) er henholdsvis tids- og mærkespecifikke effekter, og $v_{jt}$ er et fejlled. 

\begin{itemize}
    \item Estimer først first stage-regressionen \eqref{eq:first stage rxr}, og test, om \texttt{RXR} er en signifikant prædiktor for prisen ved hjælp af et t-test af \( \gamma_1 \).
    \item  Estimer herefter den strukturelle efterspørgselsligning \eqref{eq:make_fixed_effects_model} ved hjælp af 2SLS, hvor \texttt{RXR} anvendes som instrument for den endogene pris \( p_{jt} \). Sammenlign 2SLS-estimatet af \( \alpha \) med OLS-estimatet og vurder, om 2SLS-estimatet stemmer overens med din forventning om endogenitetsbias i OLS-estimatet.
    \item Argumentér for, at \texttt{RXR} er et gyldigt instrument for prisen \( p_{jt} \). Henvis til den teoretiske model.
\end{itemize}

%%%%%%%%%%%%%%%%%%%%%%%%%%%%%%%%%%%%%%%%%%%%%%%%%%%%%%%%%%%%%%%
% Structural Demand Analysis and Policy Simulation
%%%%%%%%%%%%%%%%%%%%%%%%%%%%%%%%%%%%%%%%%%%%%%%%%%%%%%%%%%%%%%%
\bigskip
\section*{Opgave 3: Gandhi Houde-instrumenterne (25\%)}

\noindent \textbf{(3A)} \cite{gandhi2019measuring} foreslår en anden tilgang til endogenitetsproblemet, nemlig at anvende instrumenter baseret på ligheder i produktkarakteristika. Ideen bag Gandhi-Houde (GH) instrumenterne er, at mere intensiv
konkurrence inden for specifikke karakteristika (f.eks. brændstofeffektivitet) reducerer producenternes markup og dermed prisen. Konkret måler et GH-instrument antallet af konkurrerende bilmodeller, der minder tilstrækkeligt meget om den undersøgte bilmodel med hensyn til en specifik karakteristik. 

Datasættet indeholder fem GH-instrumenter, hhv. \texttt{gh\_height}, \texttt{gh\_footprint}, \texttt{gh\_hp}, \texttt{gh\_mpg} og \texttt{gh\_weight}, som er konstrueret på baggrund af karakteristikaene \texttt{height}, \texttt{footprint}, \texttt{hp}, \texttt{mpg} og \texttt{weight}. 

 \textbf{OBS:} For at kunne løse denne opgave tilfredsstillende er det \textit{ikke} nødvendigt at forstå de tekniske detaljer om, hvordan GH-instrumenterne er konstrueret. Men for en god ordens skyld følger den formelle definition af GH-instrumenterne alligevel her:
\begin{equation}
    z_{jt}^{(g)} = \sum_{k \in J} 1\left(|x_{jt}^{(g)} - x_{kt}^{(g)}| < \sigma_{gt}\right),\label{eq:GH}
\end{equation}
hvor $z^{(g)}_{jt}$ er GH-instrumentet for karakteristik $g$ for bil $j$ på tidspunktet $t$, $1(\cdot)$ er indikatorfunktionen, \( x_{jt}^{(g)} \) er den normaliserede karakteristik $g$ for bil \( j \) på tidspunkt \( t \), og \( \sigma_{gt} \) er standardafvigelsen for   karakteristik $g$ beregnet for år \( t \).

\begin{itemize}
    \item Diskutér, hvorfor instrumenter som disse er relevante og forventes at være gyldige til at identificere parameteren \(\alpha\). Tag udgangspunkt i den økonometriske teori og den intuitive forklaring af GH-instrumenterne; du behøver ikke henvise til ligning \eqref{eq:GH}.
    \item Estimér efterspørgselsligningen ved hjælp af 2SLS med hhv. GH-instrumenterne alene og GH-instrumenterne sammen med RXR-instrumentet. Udfør et overidentifikationstest for at vurdere instrumenternes eksogenitet.
    \item Sammenlign resultatet med tidligere estimater og diskuter eventuelle forskelle. Overvej fordele og ulemper ved at bruge GH-instrumenterne sammenlignet med RXR.
\end{itemize} 
% \hrulefill
 \bigskip 

\noindent \textbf{(3B)} Vi skal nu bruge de estimerede efterspørgselsligninger til at beregne markedsandele, marginalomkostninger og priselasticiteter for forskellige bilmodeller over tid. Med et estimat for \(\alpha\), samt de observerede priser og markedsandele, kan vi beregne marginalomkostningerne som:
\begin{align}
    p_{jt} &= mc_{jt} + \frac{1}{- \alpha (1 - s_{jt})} \\   
    \iff mc_{jt} &= p_{jt} - \frac{1}{- \alpha (1 - s_{jt})}.
\end{align}
Egenpriselasticiteten kan beregnes som:
$$
\varepsilon_{jt} = \frac{\partial s_{jt}}{\partial p_{jt}} \cdot \frac{p_{jt}}{s_{jt}} = \alpha \cdot p_{jt} \cdot (1 - s_{jt}).
$$
Markup kan beregnes som:
$$
\text{Markup}_{jt} = \frac{p_{jt} - mc_{jt}}{p_{jt}} 
$$

\begin{itemize}
    \item  Beregn marginalomkostningerne, markups og priselasticiteter for alle observationer i datasættet baseret på både OLS og IV-estimaterne opnået i 2A og 2C, \(\hat{\alpha}^{OLS}\) og \(\hat{\alpha}^{IV}\). Har du ikke estimeret $\alpha$, så brug værdierne $\hat{\alpha}^{OLS} = -0.03$ og $\hat{\alpha}^{IV} = -0.17$.
    \item Udfør en deskriptiv analyse for at undersøge, hvordan fordelingen af marginalomkostninger, priselasticiteter og markups varierer mellem dine OLS- og IV-estimater. Illustrer også udviklingen i gennemsnitlige marginalomkostninger, markups og priselasticiteter for biltyperne \texttt{car, suv, truck} og \texttt{van} over tid på baggrund af dine IV-estimater. Brug gerne visualiseringer. 
    \item Reflektér over, hvor afhængige vi er af modeller og antagelser, når vi estimerer marginalomkostninger. Diskutér, hvilke konsekvenser det kan have, hvis modellen er fejlspecificeret, eller hvis instrumenterne ikke er valide eller stærke nok til at korrigere for prisendogenitet. 
\end{itemize}
\hrulefill

\bigskip \bigskip

\noindent \textbf{(3C)} Vi kan bruge prædikterede værdier af marginalomkostningerne fra 3B baseret på IV-estimaterne for \(\alpha\) til at estimere marginalomkostningsfunktionen \eqref{eq:mc} med bilmærke- og tidsspecifikke effekter:
\begin{equation}
    \log(mc_{jt}) = \gamma_x \mathbf{x}_{jt} + \gamma_{RXR} RXR_{jt}  + \phi_m +  \kappa_t  + \nu_{jt}
\end{equation}
hvor $\mathbf{x}_{jt}$ indeholder et konstantled og de samme produktkarakteristika, som vi anvendte i specifikation 4 i opgave 1B.
\begin{itemize}
    \item Estimer marginalomkostningsfunktionen og fortolk parameterestimaterne \(\gamma_x\) og \(\gamma_{RXR}\). Hvordan varierer marginalomkostningerne med produktkarakteristika og RXR?
    \item Er marginalomkostningerne steget eller faldet over tid, når vi kontrollerer for produktkarakteristika, RXR, og mærkespecifikke effekter?
\end{itemize} 
\bigskip \bigskip


\section*{Opgave 4: Teoretiske opgaver (25\%)}

\noindent \textbf{(4A)} Betragt følgende simple strukturelle model for udbud og efterspørgsel på bilmarkedet:
\begin{align}
y_i^D &= \beta_0 + \beta_1 p_{i} + \xi_{i}, \quad i = 1,\dots,n \label{eq:Q4 dem} \\
y_{i}^S &= \gamma_0 + \gamma_1 p_i + \gamma_2 z_{i} + \nu_{i}, \label{eq:Q4 sup}
\end{align}
hvor $y_i^D$ og $y_i^S$ repræsenterer hhv. den totale efterspørgsel og det totale udbud for bil $i$.  $p_i$ er prisen, $z_i$ er en forklarende variabel, og $\xi_i, \nu_i$ er to fejlled. Vi antager at $\beta_1 <0$ og $\gamma_1 >0$, således at efterspørgslen falder og udbuddet stiger med prisen. 

\begin{itemize}
    \item Vis, at når der er ligevægt på markedet $\left(y_i=y_i^D=y_i^S\right)$, så er den reducerede form for $p_i$:
\begin{equation}    
p_{i} = \eta_0 + \eta_1 z_{i} + u_{i}, \label{eq:Q4 p_red}
\end{equation}
hvor vi har defineret 
$$\eta_0 \equiv \frac{\gamma_0 - \beta_0}{\beta_1-\gamma_1}, \quad \eta_1 \equiv \frac{\gamma_2}{\beta_1-\gamma_1}, \quad  u_i \equiv \frac{\nu_i - \xi_i}{\beta_1-\gamma_1}.$$

Er prisen eksogen i efterspørgselsligningen, eller er den korreleret med fejlledet $\xi_i$?

\item Indsæt ligning \eqref{eq:Q4 p_red} i ligning \eqref{eq:Q4 dem} og udled den reducerede form for $y_i$ som en funktion af \(z_{i}\).

\item Diskutér med afsæt i den økonometriske teori, hvorvidt det er muligt at identificere \(\beta_1\) udelukkende baseret på den reducerede form for \(y_i\). Hvis ikke, hvad er årsagen til dette, og hvordan kan vi løse identifikationsproblemet? 
\end{itemize}
\hrulefill  
\bigskip \bigskip

\noindent \textbf{(4B)} Antag nu, at $z_i$ er et eksogent instrument, dvs.  $\text{Cov}(z, \xi) =\text{Cov}(z, \nu)= 0$. Antag desuden, at fejlledene er ukorrelerede, dvs. $\text{Cov}(\xi, \nu) =0.$

\begin{itemize}
    \item Udled biaset for OLS, når \(y_i\) 'naivt' regresseres mod prisen og en konstant. Vis at vi kan nå følgende udtryk:
    \[
 \operatorname{plim} \hat \beta^{OLS}_1 - \beta_1 =  \frac{-\,\operatorname{Var}(\xi)}{(\beta_1-\gamma_1)\,\operatorname{Var}(p)}.
\]
I hvilken retning er OLS biased? (\textit{Hint}: Husk at $\beta_1<0$ og $\gamma_1>0$ per standardantagelse). 
    \item Vis, at IV-estimatoren er konsistent for priskoefficienten \(\beta_1\).  
    \item Under antagelse af homoskedasticitet kan variansen af IV-estimatoren skrives som
\[
\text{Var}\left(\hat{\beta}^{IV}_1\right) = \frac{\operatorname{Var}{(\xi)}}{n\operatorname{Cov}(z, p)^2},
\] 
Vis, at vi i vores model kan udlede følgende udtryk for IV-estimatorens varians:
\[
\text{Var}\left(\hat{\beta}^{IV}_1\right) = \frac{\operatorname{Var}(\xi)}{n\eta_1^2 \operatorname{Var}(z)^2} = \frac{(\beta_1-\gamma_1)^2 \operatorname{Var}(\xi)}{n \gamma_2^2 \operatorname{Var}(z)^2}.
\]
Hvordan påvirker styrken af instrumentet \(z\) IV-estimatorens konsistens og varians? Hvordan bliver variansen påvirket af hældningen på udbuds- og efterspørgselskurverne  $\beta_1$ og $\gamma_1$? Fortolk.
\end{itemize}
\hrulefill

\bigskip \bigskip
\noindent \textbf{(4C)} Lad os undersøge de teoretiske forudsigelser fra opgave 4B gennem et Monte Carlo eksperiment. 
\begin{itemize}
    \item Simuler data med $n= 500$ observationer og $r=500$ replikationer baseret på ligevægtsmodellen givet ved \eqref{eq:Q4 dem} og \eqref{eq:Q4 p_red}, dvs. følgende ligninger:
    \begin{align*}
        y_i &= \beta_0 + \beta_1 p_{i} + \xi_{i} \\
        p_{i} &= \eta_0 + \eta_1 z_{i} + u_{i}, \label{eq:Q4 p_red}
    \end{align*}
    hvor vi har defineret
    $$\eta_0 \equiv \frac{\gamma_0 - \beta_0}{\beta_1-\gamma_1}, \quad \eta_1 \equiv \frac{\gamma_2}{\beta_1-\gamma_1}, \quad  u_i \equiv \frac{\nu_i - \xi_i}{\beta_1-\gamma_1}.
    $$
    Brug følgende baseline parameterværdier:
 \(\beta_0 = 0.5\), $\beta_1=-0.1$ \(\gamma_0 = 0.25\), \(\gamma_1 = 1\), $\gamma_2=1$ og antag, at \(z_{i} \sim N(0,1)\), \(\xi_{i} \sim N(0,1)\), og \(\nu_{i} \sim N(0,1)\). 
    
    Estimer for hver replikation \(\beta_1\) ved OLS og ved IV, hvor \(z_{i}\) bruges som instrument for \(p_{i}\), og beregn bias og varians af OLS- og IV-estimatorerne på tværs af replikationerne. 
    \item Lad os kalde ovenstående scenarie 1. Gentag nu eksperimentet i følgende supplerende scenarier:
    \begin{itemize}
      \item Scenarie 2: Brug baseline parametrene, men sæt $\gamma_2 = 0.25$
      \item Scenarie 3: Brug baseline parametrene, men sæt $\gamma_1 = 0.25$ 
      \end{itemize}
      Hvad er den økonomiske fortolkning af forskellene mellem de tre scenariers parameterværdier?
    \item  Diskutér simulationsresultaterne med fokus på, hvordan parameterværdierne påvirker estimatorernes varians og bias på tværs af de tre scenarier.
\end{itemize}





\break

\appendix
\section*{Appendix}
\subsection*{Model for udbud, efterspørgsel og ligevægt på bilmarkedet (til opgave 1-3)}
\setcounter{equation}{0}\renewcommand\theequation{A\arabic{equation}}

\noindent\emph{Forbrugernes efterspørgsel efter biler}\\
Vi starter med at modellere forbrugernes valg af biler. Forbrugerne maksimerer deres nytte ved valg mellem forskellige bilmodeller, 
$j=0,1,\dots,J$. Vi antager, at forbrugerne vælger det produkt, $j^*$, der maksimerer deres nytte, givet prisen og produktkarakteristika for produkt $j$ på tidspunkt $t$:
\[
j_{it}^* = \arg \max_{j=0,1,\dots, J} U_{ijt},
\]
hvor $j=0$ repræsenterer et outside good (dvs. valget at ikke købe en bil).

Vi antager, at forbrugernes præferencer kan beskrives ved en lineær nyttefunktion
\[
U_{ijt} = \beta \mathbf{x}_{jt} + \alpha p_{jt} + \xi_{jt} + \epsilon_{ijt},
\]
hvor $\mathbf{x}_{jt}$ er observerbare karakteristika for produkt $j$ på tidspunkt $t$ (f.eks. hestekræfter, brændstofeffektivitet), $p_{jt}$ er prisen på produkt $j$ på tidspunkt $t$, $\xi_{jt}$ er uobserverbare karakteristika, som kan være brand-image eller andre kvaliteter, der ikke er direkte målt, og $\epsilon_{ijt}$ er den idiosynkratiske præference for forbruger $i$, som antages at følge en ekstremværdifordeling. 

Antagelsen om, at $\epsilon_{ijt}$ følger en ekstremværdifordeling, leder til aggregerede markedsandele på logit-formen
\begin{equation*}
s_{jt} = \frac{\exp(\beta \mathbf{x}_{jt} + \alpha p_{jt} + \xi_{jt})}{1 + \sum_{k=1}^J \exp(\beta \mathbf{x}_{kt} + \alpha p_{kt} + \xi_{kt})},
\vspace{1em}
\end{equation*}
hvor vi har normaliseret nytten af vores outside good til 0, og hvor $s_{jt}$ er markedsandelen for produkt $j$ på tidspunkt $t$.

Vi kan skrive log-markedsandelen relativt til andelen af husholdninger uden bil:
\begin{equation}
    \log\left(\frac{s_{jt}}{s_{0t}}\right) = \beta \mathbf{x}_{jt} + \alpha p_{jt} + \xi_{jt}. \label{eq:A-demand}    
\end{equation}
hvilket resulterer i en lineær regressionsmodel, hvor $\log\left(\frac{s_{jt}}{s_{0t}}\right)$ er den afhængige variabel, $\mathbf{x}_{jt}$ og $p_{jt}$ er observerbare forklarende variable, $\xi_{jt}$ er fejlleddet, og $\beta$ og $\alpha$ er parametre vi ønsker at estimere.
\\ \medskip 

\noindent\emph{Virksomhedernes prissætning og marginalomkostninger}\\
Vi antager, at virksomhederne sætter prisen $p_{jt}$, der maksimerer profit givet efterspørgselsfunktionen for produkt $j$ på tidspunktet $t$
$$ p_{jt} = \arg \max_{p} \left[(p - mc_{jt}) s_{jt} \right],$$
hvor $mc_{jt}$ er marginalomkostningerne for produkt $j$ på tidspunktet $t$. Førsteordensbetingelsen for optimal prissætning er givet ved
\[
s_{jt} + (p_{jt} - mc_{jt}) \frac{\partial s_{jt}}{\partial p_{jt}} = 0.
\]
Denne betingelse giver den optimale prissætning som en funktion af marginalomkostningerne og efterspørgslen:
\begin{align}
    p_{jt} = mc_{jt} - s_{jt}\left(\frac{\partial s_{jt}}{\partial p_{jt}}\right)^{-1} = mc_{jt} + \frac{1}{- \alpha (1 - s_{jt})}. \label{eq:A-supply}    
\end{align}
hvor det sidste lighedstegn følger ved at differentiere efterspørgselsfunktionen i \eqref{eq:A-demand} mht. til prisen $\frac{\partial s_{jt}}{\partial p_{jt}}={\alpha s_{jt}(1 - s_{jt})}$. 

Vi antager, at logaritmen til marginalomkostningerne for virksomhed $j$ i marked $t$ kan beskrives som en lineær funktion af produktkarakteristika, $\mathbf{x}_{jt}$ og udbudsspecifikke karakteristika, $\mathbf{w}_{jt}$:
\begin{equation}
    \log(mc_{jt}) = \gamma_x \mathbf{x}_{jt} + \gamma_w \mathbf{w}_{jt} + \nu_{jt}, \label{eq:A-mc}    
\end{equation}
hvor $\mathbf{x}_{jt}$ er de observerbare produktkarakteristika som også påvirker nytten (f.eks. størrelse, kvalitet), $\mathbf{w}_{jt}$ er udbudsspecifikke karakteristika (f.eks. råvarepriser, skatter, toldmure, lønomkostninger, mv.), $\nu_{jt}$ er et stokastisk udbudsstød, og $\gamma_x$ og $\gamma_w$ er parametre, vi ønsker at estimere. 

Således kan modellen opsummeres ved ligningerne \eqref{eq:A-demand}, \eqref{eq:A-supply} og \eqref{eq:A-mc}.

\vspace{1em}
\noindent\hrulefill

\bibliography{references}

\end{document}

